% =============================================================================
% Kapitel 3: Systementwurf und Implementierung
% =============================================================================

\chapter{Systementwurf und Implementierung}
\label{chap:implementierung}

Dieses Kapitel beschreibt den Entwurf und die prototypische Umsetzung des
Off"=Grid"=Schutzsensorsystems. Ausgehend von den in Kapitel~\ref{chap:einleitung}
abgeleiteten Anforderungen werden zunächst die gewählten Hardware"=Komponenten
vorgestellt und begründet. Anschließend wird die Systemarchitektur auf Software"=Ebene
beschrieben, bevor auf die konkrete Implementierung der Firmware, des Network"=Servers
und der Datenvisualisierung eingegangen wird.

% TODO: Blockschaltbild des Gesamtsystems einfügen (Sensoreinheit → Gateway → RPi Network-Server)

% -----------------------------------------------------------------------------
\section{Hardwareauswahl}
\label{sec:hardwareauswahl}
% -----------------------------------------------------------------------------

\subsection{Mikrocontroller}
\label{subsec:mcu}

\paragraph{Mikrocontroller}

Für die Sensoreinheit wird der \textbf{STM32WL55CC} von STMicroelectronics
eingesetzt, der auf dem Entwicklungsboard \textbf{NUCLEO"=WL55JC1} verfügbar
ist. Der STM32WL55 ist der erste am Markt verfügbare \acrfull{soc},
der einen ARM Cortex"=M4"=Applikationsprozessor, einen ARM Cortex"=M0+ für
den LoRa"=Stack sowie einen integrierten Sub"=GHz"=Transceiver (SX126x"=kompatibel)
auf einem einzigen Die vereint. Dies eliminiert den Bedarf an einem separaten
Funkmodul und reduziert damit Platzbedarf, Kostenpunkt und potenzielle
Fehlerquellen an der Schnittstelle zwischen MCU und Funkmodul.

Wichtige Eigenschaften des STM32WL55 für dieses Projekt lassen sich dem Datenblatt entnehmen:

\begin{itemize}
    \item \textbf{Integrierter Sub"=GHz"=Transceiver:} Unterstützt LoRa (CSS),
    \acrshort{fsk} und \acrshort{bpsk} im Frequenzbereich 150--960\,MHz,
    damit nativ kompatibel mit dem europäischen \SI{868}{\MHz}"=Band.

    \item \textbf{Sendeleistung:} bis zu \SI{22}{dBm} bei
    $V_\text{DDPA} = \SI{3{,}3}{\V}$, was -- in Kombination mit SF12 --
    Reichweiten von mehreren Kilometern im Freifeld ermöglicht.

    \item \textbf{Energieprofil:} Im Deep"=Stop"=Modus beträgt die
    Stromaufnahme lediglich \SI{1{,}08}{\micro\ampere}, was einen
    mehrschichtigen Betrieb auf Primärzellen ermöglicht.

    \item \textbf{Zephyr-Support:} Das Board \texttt{nucleo\_wl55jc}
    ist offiziell von Zephyr-\acrshort{rtos} unterstützt, was einen produktionsreifen
    LoRaWAN"=Stack ohne proprietäre Middleware ermöglicht.

    \item \textbf{Integrierter ST"=LINK V3E:} Das NUCLEO"=Board enthält
    einen integrierten Debugger und Programmer, sodass kein separates
    Programmiergerät erforderlich ist.
\end{itemize}

\paragraph{Gateway}
% TODO: write sth about DIY RPi as Gateway

TODO aaaa

\paragraph{Messinstrument}

Der \textbf{Nordic Power Profiler Kit~II (PPK2)} ermöglicht eine
hochauflösende Strommessung im Bereich von \SI{100}{\nano\ampere} bis
\SI{1}{\ampere} und dient damit der Verifikation der
theoretischen Energiebudgetberechnungen. Er wird zwischen Spannungsquelle
und NUCLEO"=Board geschaltet und zeichnet das Stromprofil über die Zeit auf,
sodass Schlafphasen, Aufwachvorgänge und Sendepulse direkt sichtbar werden.

\paragraph{Radioempfänger}

Der \textbf{RTL-SDR V4} ist ein kostengünstiger \acrfull{sdr}"=Empfänger,
der in Kombination mit SDR\# (Windows) oder GNU Radio (Linux) zur Spektrumanalyse
im \SI{868}{\MHz}"=Band eingesetzt wird. LoRa"=Chirps sind im Wasserfall"=Diagramm
als charakteristische diagonale Linien sichtbar. Mit dem GNU Radio"=Plugin
\texttt{gr-lora} können LoRaWAN"=Pakete vollständig demoduliert und dekodiert
werden, ohne dass ein zweites LoRa"=Endgerät als Empfänger erforderlich ist.

% -----------------------------------------------------------------------------
\section{Software-Architektur}
\label{sec:softwarearchitektur}
% -----------------------------------------------------------------------------

\subsection{Überblick der Softwareschichten}
\label{subsec:sw-überblick}

Die Softwarearchitektur des Systems gliedert sich in drei Ebenen, die in
Abbildung~\ref{fig:sw_stack} dargestellt sind: % TODO: Abbildung erstellen

\begin{itemize}
    \item \textbf{Firmware-Ebene (Endgerät):} Zephyr RTOS mit integriertem
    LoRaWAN"=Stack auf dem STM32WL55. Zuständig für Sensorauslese,
    Schlafmodussteuerung und Paketübertragung.

    \item \textbf{Netzwerkebene (Gateway + Network"=Server):} ChirpStack
    als lokaler, offline"=fähiger LoRaWAN"=Network"=Server auf einem
    Raspberry Pi. Zuständig für Paketempfang, Authentifizierung,
    ADR"=Steuerung und Weiterleitung an die Applikationsebene.

    \item \textbf{Applikationsebene:} Node"=RED zur visuellen Datenpipeline
    und Echtzeit"=Visualisierung der Messwerte. MQTT als Kommunikationsprotokoll
    zwischen ChirpStack und Node"=RED, gleichzeitig als Debugging"=Schnittstelle
    nutzbar über den MQTT Explorer.
\end{itemize}

\subsection{Auswahl des Network-Servers: ChirpStack}
\label{subsec:chirpstack}

Für den Network"=Server"=Betrieb wurde \textbf{ChirpStack} gegenüber The Things
Stack (TTS) gewählt. Der entscheidende Grund ist die vollständige
Offline"=Fähigkeit: ChirpStack kann auf einem Raspberry Pi im lokalen Netz
betrieben werden, ohne dass eine Verbindung zu externen Cloud"=Servern
erforderlich ist. TTS Community Edition setzt hingegen eine permanente
Internetverbindung zu den TTN"=Servern voraus -- eine Anforderung, die im
Off"=Grid"=Szenario strukturell nicht erfüllbar ist. Für die frühe
Entwicklungsphase und erste Funktionstests kann TTS dennoch als schnell
verfügbare Cloud"=Infrastruktur genutzt werden, bevor das lokale
ChirpStack"=Setup vollständig konfiguriert ist.

% -----------------------------------------------------------------------------
\section{Firmware-Implementierung}
\label{sec:firmware}
% -----------------------------------------------------------------------------

% TODO: Diesen Abschnitt ausarbeiten sobald Implementierung abgeschlossen

\subsection{Zephyr RTOS und Projektstruktur}
% Devicetree-Overlay für Gassensor, prj.conf Kconfig-Flags,
% West build/flash Workflow, App-Struktur

\subsection{LoRaWAN-Stack-Konfiguration}
% OTAA Join-Prozedur, Kconfig-Flags (CONFIG_LORAWAN=y etc.),
% Payload-Kodierung, Frame Counter Persistenz

\subsection{Power-Management-Implementierung}
% Zephyr PM Framework, Deep-Stop-Modus STM32WL55,
% Tickless Idle, PM_DEVICE_STATE_SUSPEND für Peripherie,
% Stromprofil: aktive Phase / TX-Phase / Sleep-Phase

\subsection{Simulation mit Renode}
% STM32WL55 Platform-File (.repl), .resc Testskript,
% Firmware-Test ohne Hardware, west simulate --runner=renode


% -----------------------------------------------------------------------------
\section{Gateway- und Network-Server-Konfiguration}
\label{sec:gateway-ns}
% -----------------------------------------------------------------------------

% TODO: Diesen Abschnitt ausarbeiten sobald Hardware entschieden

\subsection{Gateway-Inbetriebnahme}
% Paketweiterleiter-Konfiguration (UDP Packet Forwarder / TS008),
% Kanal-Plan für EU868, Gateway-EUI

\subsection{ChirpStack-Konfiguration}
% Installation auf Raspberry Pi (Docker),
% Device-Profile (OTAA, SF-Konfiguration),
% Application anlegen, MQTT-Integration


% -----------------------------------------------------------------------------
\section{Datenvisualisierung}
\label{sec:visualisierung}
% -----------------------------------------------------------------------------

% TODO: Diesen Abschnitt ausarbeiten sobald Pipeline steht

\subsection{MQTT-Pipeline}
% ChirpStack → MQTT Broker → Node-RED → Dashboard
% Relevante Topics: application/+/device/+/event/up

\subsection{Kartografische Auswertung der Messdaten}
% GPS-Koordinaten + RSSI als CSV exportieren,
% Python + Folium: Heatmap über OpenStreetMap,
% Overlay mit SPLAT! Coverage-Karte für Theorie-Praxis-Vergleich

\newacronym{soc}{SoC}{System-on-Chip}
\newacronym{bpsk}{BPSK}{Binary Phase Shift Keying}
\newacronym{sdr}{SDR}{Software Defined Radio}
\newacronym{rtos}{RTOS}{Real-Time Operating System}